\section{Deterministic one-step estimates}
\label{app:det-proofs}

We prove the two one-step estimates used in \Cref{sec:global}, namely
\Cref{lem:descent-Riem} for the Riemannian retraction and \Cref{prop:onestep-KL} for the
KL/Bregman step, together with the relative smoothness of $\calE_2$ that identifies the
Bregman geometry of the latter. Throughout $a=(m,C)$ is a Gaussian parameter, $g:=g(a)$,
$A:=A(a)$, $R:=R(a)=C^{-1}-A$, and $\calE=\calE_1+\calE_2$ as in \eqref{eq:E-split}. Only
convexity and $\beta$-smoothness of $V$, that is $0\preceq\nabla^2V\preceq\beta\Id$, are used;
both hold under \Cref{assump:logconcave-smooth}.

\subsection{Two scalar inequalities}

\begin{lemma}\label{lem:scalar-retraction}
    For every $w\in[-1,1]$,
    \begin{equation*}
        \left(e^{w/2}-1-\frac w2\right)+\frac12\left(e^{w/2}-1\right)^2
        \leq\frac{w^2}2 .
    \end{equation*}
\end{lemma}

\begin{proof}
    Put $v=w/2\in[-1/2,1/2]$. Since
    $(e^v-1)+\tfrac12(e^v-1)^2=\tfrac12(e^{2v}-1)$, the desired inequality is equivalent, with
    $x=2v=w\in[-1,1]$, to $e^x\leq1+x+x^2$. Indeed
    \begin{equation*}
        e^x-1-x=\sum_{k\geq2}\frac{x^k}{k!}
        \leq x^2\sum_{k\geq2}\frac{|x|^{k-2}}{k!}
        \leq(e-2)x^2<x^2,
    \end{equation*}
    which proves the claim.
\end{proof}

\begin{lemma}\label{lem:scalar-KL}
    Let $K\geq1$, $0<\Delta t\leq1/K$, $\mu\in[0,K]$, and
    $b:=\sqrt{(1+\Delta t)/(1+\Delta t\,\mu)}$. Then
    \begin{equation}\label{eq:scalar-KL}
        \mu(b-1)-\log b+\frac K2(b-1)^2
        \leq
        -\frac1{2\Delta t}\bigl(b^2-1-2\log b\bigr).
    \end{equation}
\end{lemma}

\begin{proof}
    Because $K\leq1/\Delta t$ and $(b-1)^2\geq0$, it suffices to replace $K$ on the left by
    $1/\Delta t$. The relation between $b$ and $\mu$ is
    $\mu=((1+\Delta t)b^{-2}-1)/\Delta t$. Multiplying the resulting inequality by $\Delta t$
    and moving all terms to the left, define
    \begin{equation*}
        G(b):=\Delta t\,\mu(b-1)-\Delta t\log b+\tfrac12(b-1)^2+\tfrac12(b^2-1-2\log b),
    \end{equation*}
    so that $G(1)=0$ and a direct differentiation gives
    \begin{equation*}
        G'(b)=\frac{(b-1)\bigl(2b^3-(1+\Delta t)(b+2)\bigr)}{b^3}.
    \end{equation*}
    Since $0\leq\Delta t\,\mu\leq\Delta t\,K\leq1$, we have
    $b\in[\sqrt{(1+\Delta t)/2},\sqrt{1+\Delta t}]$. The polynomial
    $p(b):=2b^3-(1+\Delta t)(b+2)$ is convex on $(0,\infty)$, and at the two endpoints of this
    interval
    \begin{equation*}
        p\bigl(\sqrt{(1+\Delta t)/2}\bigr)=-2(1+\Delta t)<0,
        \qquad
        p\bigl(\sqrt{1+\Delta t}\bigr)=(1+\Delta t)\bigl(\sqrt{1+\Delta t}-2\bigr)<0,
    \end{equation*}
    using $\Delta t\leq1$. Convexity therefore implies $p(b)<0$ throughout the interval, so
    $G'(b)\geq0$ for $b\leq1$ and $G'(b)\leq0$ for $b\geq1$. Thus $G$ attains its maximum at
    $b=1$, where it vanishes, proving \eqref{eq:scalar-KL}.
\end{proof}

\subsection{Retraction smoothness and the Riemannian descent estimate}

\begin{proof}[Proof of \Cref{lem:descent-Riem}]
    Let $Z\sim\N(0,\Id)$ and introduce the coupling
    \begin{equation*}
        X:=m+C^{1/2}Z\sim\rho_a,
        \qquad
        Y:=m+u+C^{1/2}e^{W/2}Z\sim\rho_{\Retr_a(\zeta)},
    \end{equation*}
    which is legitimate because $e^{W/2}$ is symmetric, so $Y$ is Gaussian with mean $m+u$ and
    covariance $C^{1/2}e^WC^{1/2}$. Write $E_W:=e^{W/2}-\Id$ and $M:=C^{1/2}AC^{1/2}$. The
    entropy change is exact,
    \begin{equation*}
        \calE_1(\Retr_a(\zeta))-\calE_1(a)=-\tfrac12\Tr W,
    \end{equation*}
    and by $\beta$-smoothness of $V$,
    \begin{equation*}
        \calE_2(\Retr_a(\zeta))-\calE_2(a)
        \leq
        \E\bigl[\nabla V(X)^\top(Y-X)\bigr]+\frac\beta2\E\norm{Y-X}^2 .
    \end{equation*}
    Because $Y-X=u+C^{1/2}E_WZ$, $\E[\nabla V(X)]=-g$, and Gaussian integration by parts gives
    $\E[Z\nabla V(X)^\top]=C^{1/2}A$, the first-order contribution equals
    $-g^\top u+\Tr(E_WM)$. The cross term in the quadratic contribution vanishes, and
    $C\preceq\Lambda\Id$ yields
    \begin{equation*}
        \frac\beta2\E\norm{Y-X}^2
        \leq
        \frac K2u^\top C^{-1}u+\frac K2\norm{E_W}_\Frob^2 .
    \end{equation*}
    On the other hand, since $\grad\calE(a)=(-Cg,-CRC)$ and $C^{1/2}RC^{1/2}=\Id-M$,
    \begin{equation*}
        \inner{\grad\calE(a)}\zeta_a=-g^\top u-\tfrac12\Tr((\Id-M)W).
    \end{equation*}
    Subtracting this linear term from the entropy and covariance contributions leaves
    \begin{equation*}
        \Theta
        :=\Tr\left[\left(e^{W/2}-\Id-\frac W2\right)M\right]
          +\frac K2\norm{e^{W/2}-\Id}_\Frob^2 .
    \end{equation*}
    The matrix $N_W:=e^{W/2}-\Id-W/2$ is positive semidefinite while $0\preceq M\preceq K\Id$,
    hence $\Tr(N_WM)\leq K\Tr N_W$. Diagonalizing $W$ and applying
    \Cref{lem:scalar-retraction} eigenvalue by eigenvalue gives
    \begin{equation*}
        \Theta
        \leq K\sum_i\left[e^{w_i/2}-1-\frac{w_i}2+\frac12(e^{w_i/2}-1)^2\right]
        \leq\frac K2\norm W_\Frob^2 .
    \end{equation*}
    Combining the mean and covariance estimates and using
    $\norm\zeta_a^2=u^\top C^{-1}u+\tfrac12\norm W_\Frob^2$ proves
    \eqref{eq:Riem-retraction-smooth}.

    For the deterministic step take $\zeta=-\Delta t\grad\calE(a)=(\Delta t\,Cg,\Delta t\,CRC)$,
    so that
    \begin{equation*}
        W=\Delta t\,C^{1/2}RC^{1/2}=\Delta t\bigl(\Id-C^{1/2}AC^{1/2}\bigr).
    \end{equation*}
    The eigenvalues of the matrix in parentheses lie in $[1-K,1]$, so
    $\norm W_\op\leq\Delta t\,K\leq1$. Moreover
    $\inner{\grad\calE(a)}\zeta_a=-\Delta t\,\Gradnorm(a)^2$ and
    $\norm\zeta_a^2=\Delta t^2\Gradnorm(a)^2$, and substitution into
    \eqref{eq:Riem-retraction-smooth} gives
    \begin{equation*}
        \calE(a_{n+1})\leq\calE(a_n)-\Delta t(1-K\Delta t)\Gradnorm(a_n)^2 .
    \end{equation*}
    The final assertion is the case $\Delta t\leq1/(2K)$.
\end{proof}

\subsection{The KL/Bregman step: variational characterization}

\begin{lemma}\label{lem:KL-argmin}
    The update \eqref{upd:KL} is the forward--backward proximal step
    \begin{equation}\label{eq:KL-argmin}
        a_{n+1}
        =\argmin_{a'=(m',C')}
        \Bigl\{
            \Delta t\inner{\nabla\calE_2(a_n)}{a'-a_n}
            +\Delta t\,\calE_1(a')
            +\KL(\rho_{a'}\|\rho_{a_n})
        \Bigr\},
    \end{equation}
    which linearizes the potential $\calE_2$ explicitly and treats the entropy $\calE_1$
    implicitly, and its minimizer is
    \begin{equation*}
        m_{n+1}=m_n+\Delta t\,C_ng(a_n),
        \qquad
        C_{n+1}=(1+\Delta t)\bigl(C_n^{-1}+\Delta t\,A(a_n)\bigr)^{-1}.
    \end{equation*}
\end{lemma}

\begin{proof}
    The Gaussian KL splits as
    \begin{equation*}
        \KL(\rho_{a'}\|\rho_{a_n})
        =\tfrac12(m_n-m')^\top C_n^{-1}(m_n-m')
        +\tfrac12\bigl[\Tr(C_n^{-1}C')-N_\theta+\log\det C_n-\log\det C'\bigr],
    \end{equation*}
    so the objective in \eqref{eq:KL-argmin} decouples in $m'$ and $C'$. The $m'$-part is
    $-\Delta t\,g(a_n)^\top(m'-m_n)+\tfrac12(m'-m_n)^\top C_n^{-1}(m'-m_n)$, minimized at
    $m'-m_n=\Delta t\,C_ng(a_n)$. The $C'$-part, dropping $C'$-independent terms and using
    $\nabla_C\calE_2=\tfrac12A$ and $\calE_1(C')=-\tfrac12\log\det C'$, is
    \begin{equation*}
        \tfrac{\Delta t}2\Tr(A(a_n)C')+\tfrac12\Tr(C_n^{-1}C')
        -\tfrac{1+\Delta t}2\log\det C'
        =\tfrac12\Tr\bigl((C_n^{-1}+\Delta t\,A(a_n))C'\bigr)
        -\tfrac{1+\Delta t}2\log\det C',
    \end{equation*}
    whose stationarity condition
    $\tfrac12(C_n^{-1}+\Delta t\,A(a_n))=\tfrac{1+\Delta t}2(C')^{-1}$ gives the stated
    covariance. Both objectives are strictly convex in their arguments, since $C_n^{-1}\succ0$
    and $-\log\det$ is convex, so the stationary points are the minimizers.
\end{proof}

\subsection{Relative smoothness of the potential part}

\begin{proposition}\label{prop:smooth-Bregman}
    Under \Cref{assump:logconcave-smooth}, let $a_n=(m_n,C_n)$ satisfy
    $\lambda_{\min}\Id\preceq C_n\preceq\lambda_{\max}\Id$. Then $\calE_2$ is $L$-smooth with
    respect to the local Bregman divergence $D(\cdot,a_n)=\KL(\rho_\cdot\|\rho_{a_n})$ of
    \eqref{eq:Bregman}: for every $a=(m,C)$ with
    $\lambda_{\min}\Id\preceq C\preceq\lambda_{\max}\Id$,
    \begin{equation*}
        \calE_2(a)\leq\calE_2(a_n)+\inner{\nabla\calE_2(a_n)}{a-a_n}+L\,D(a,a_n),
        \qquad
        L=\beta\lambda_{\max}.
    \end{equation*}
\end{proposition}

\begin{proof}
    We consider the displacement from $a_n=(m_n,C_n)$ to $a=(m,C)$ and denote the direction
    $u=m-m_n$ and $H=B-B_n$, where $M:=C_n^{-1/2}CC_n^{-1/2}$, $B_n:=C_n^{1/2}$ and
    $B:=C_n^{1/2}M^{1/2}$, so that $BB^\top=C$. Under this formulation we write
    $\calE_2(m,B):=\E_{\xi\sim\N(0,\Id)}[V(m+B\xi)]$ and denote the trajectory
    \begin{equation*}
        \varphi(t)=\E_{\xi\sim\N(0,\Id)}\bigl[V(m_n+tu+(B_n+tH)\xi)\bigr].
    \end{equation*}
    Writing $X_t=m_n+tu+(B_n+tH)\xi$,
    \begin{equation*}
        \varphi'(t)=\E\bigl[\nabla V(X_t)^\top(u+H\xi)\bigr],
        \qquad
        \varphi''(t)=\E\bigl[(u+H\xi)^\top\nabla^2V(X_t)(u+H\xi)\bigr],
    \end{equation*}
    and $\nabla^2V\preceq\beta\Id$ gives $\varphi''(t)\leq\beta(\norm u^2+\norm H_\Frob^2)$. As
    Taylor's expansion gives $\varphi(1)=\varphi(0)+\varphi'(0)+\int_0^1(1-s)\varphi''(s)\dd s$,
    \begin{equation*}
        \varphi(1)-\varphi(0)-\varphi'(0)\leq\frac\beta2\bigl(\norm u^2+\norm H_\Frob^2\bigr),
    \end{equation*}
    which is equivalent to
    \begin{equation*}
        \calE_2(a)-\calE_2(a_n)
        -\inner{\nabla_m\calE_2(a_n)}u
        -\inner{\nabla_B\calE_2(a_n)}H
        \leq\frac\beta2\bigl(\norm u^2+\norm H_\Frob^2\bigr).
    \end{equation*}
    Denote the gradient of $\calE_2$ with respect to $C$ by
    \begin{equation*}
        G_n:=\nabla_C\calE_2(a_n)=\tfrac12\E_{\rho_{a_n}}[\nabla^2V(\theta)]=\tfrac12A(a_n),
    \end{equation*}
    where $G_n\succeq0$ because $V$ is convex. Since
    $\nabla_B\calE_2(a_n)=\nabla_C\calE_2(a_n)\cdot2B_n=2G_nB_n$,
    \begin{align*}
        \inner{\nabla_B\calE_2(a_n)}H
        &=\Tr\bigl((2G_nB_n)^\top H\bigr)=\Tr\bigl(G_n(B_nH+HB_n)\bigr)\\
        &\leq\Tr\bigl(G_n(B_nH^\top+HB_n+HH^\top)\bigr)
        =\inner{\nabla_C\calE_2(a_n)}{C-C_n},
    \end{align*}
    the inequality coming from $G_n\succeq0$ and $HH^\top\succeq0$. We therefore obtain
    \begin{equation*}
        \calE_2(a)
        \leq\calE_2(a_n)+\inner{\nabla\calE_2(a_n)}{a-a_n}
        +\frac\beta2\bigl(\norm u^2+\norm H_\Frob^2\bigr).
    \end{equation*}

    It remains to connect $\tfrac\beta2(\norm u^2+\norm H_\Frob^2)$ to the divergence $D$. For
    the mean part,
    \begin{equation*}
        \frac\beta2\norm u^2
        \leq\beta\lambda_{\max}\cdot\frac12u^\top C_n^{-1}u
        \leq\beta\lambda_{\max}D^m(m,m_n).
    \end{equation*}
    For the covariance part, $D^C(C,C_n)=\tfrac12(\Tr M-\log\det M-N_\theta)$, and the
    eigenvalues $\eta_1,\ldots,\eta_{N_\theta}$ of $M$ satisfy
    $\eta_i\in[\lambda_{\min}/\lambda_{\max},\lambda_{\max}/\lambda_{\min}]$. This gives
    \begin{equation*}
        D^C(C,C_n)
        =\frac12\sum_{i=1}^{N_\theta}(\eta_i-\log\eta_i-1)
        \geq\frac12\sum_{i=1}^{N_\theta}(\sqrt{\eta_i}-1)^2
        =\frac12\norm{M^{1/2}-\Id}_\Frob^2
        \geq\frac1{2\lambda_{\max}}\norm H_\Frob^2,
    \end{equation*}
    where the first inequality comes from $\log x\leq x-1$ and the second from
    $H=B-B_n=C_n^{1/2}(M^{1/2}-\Id)$ together with $C_n\preceq\lambda_{\max}\Id$. These
    collectively give
    $\tfrac\beta2(\norm u^2+\norm H_\Frob^2)\leq L\,D(a,a_n)$ with $L=\beta\lambda_{\max}$,
    which completes the proof.
\end{proof}

\subsection{Forward-KL descent and the Wasserstein-gradient lower bound}

\begin{proof}[Proof of \Cref{prop:onestep-KL}]
    Set
    \begin{equation*}
        M:=C_n^{1/2}A(a_n)C_n^{1/2},
        \qquad
        B:=\sqrt{1+\Delta t}\,(\Id+\Delta t\,M)^{-1/2},
    \end{equation*}
    so that $0\preceq M\preceq K\Id$, $B$ is a symmetric function of $M$, and
    $C_{n+1}=C_n^{1/2}B^2C_n^{1/2}$ by \Cref{lem:KL-argmin}. With $\xi\sim\N(0,\Id)$, couple
    \begin{equation*}
        X:=m_n+C_n^{1/2}\xi\sim\rho_{a_n},
        \qquad
        Y:=m_n+\Delta t\,C_ng+C_n^{1/2}B\xi\sim\rho_{a_{n+1}} .
    \end{equation*}
    The entropy change is $\calE_1(a_{n+1})-\calE_1(a_n)=-\log\det B$. Writing
    $E:=B-\Id$, the $\beta$-smoothness estimate of the proof of \Cref{lem:descent-Riem} applies
    verbatim with $E$ in place of $e^{W/2}-\Id$, since both are symmetric functions of $M$ and
    the computation uses only
    $V(Y)\leq V(X)+\nabla V(X)^\top(Y-X)+\tfrac\beta2\norm{Y-X}^2$, $\E[\nabla V(X)]=-g$, the
    Price identity $\E[\xi\nabla V(X)^\top]=C_n^{1/2}A$, and $C_n\preceq\Lambda\Id$. It gives
    \begin{equation*}
        \calE_2(a_{n+1})-\calE_2(a_n)
        \leq
        -\Delta t\,g^\top C_ng+\Tr(EM)
        +\frac{K\Delta t^2}2g^\top C_ng+\frac K2\norm E_\Frob^2 .
    \end{equation*}
    Because $K\Delta t\leq1$, the mean contribution is at most
    $-\tfrac{\Delta t}2g^\top C_ng$, which is exactly $-\Delta t^{-1}$ times the mean part
    $\tfrac{\Delta t^2}2g^\top C_ng$ of $\KL(\rho_{a_{n+1}}\|\rho_{a_n})$. For the covariance
    part, let $\mu_i$ be the eigenvalues of $M$ and
    $b_i=\sqrt{(1+\Delta t)/(1+\Delta t\,\mu_i)}$ the aligned eigenvalues of $B$; adding the
    entropy change, the covariance contribution is
    \begin{equation*}
        \sum_i\left[\mu_i(b_i-1)-\log b_i+\frac K2(b_i-1)^2\right],
    \end{equation*}
    which by \Cref{lem:scalar-KL} is at most
    $-\tfrac1{2\Delta t}\sum_i(b_i^2-1-2\log b_i)$. Since $C_n^{-1}C_{n+1}$ has eigenvalues
    $b_i^2$, this is exactly $-\Delta t^{-1}$ times the covariance part
    $\tfrac12\sum_i(b_i^2-1-\log b_i^2)$ of $\KL(\rho_{a_{n+1}}\|\rho_{a_n})$. Adding the mean
    and covariance parts proves the descent inequality.

    For the lower bound, the mean part of the forward divergence is
    \begin{equation*}
        \frac12(m_{n+1}-m_n)^\top C_n^{-1}(m_{n+1}-m_n)
        =\frac{\Delta t^2}2g^\top C_ng
        \geq\frac{L\Delta t^2}2\norm g^2 .
    \end{equation*}
    For the covariance part, set $S:=C_n^{1/2}RC_n^{1/2}=\Id-M$ and
    \begin{equation*}
        Q:=C_n^{1/2}C_{n+1}^{-1}C_n^{1/2}=\frac{\Id+\Delta t\,M}{1+\Delta t}.
    \end{equation*}
    For $q>0$ one has
    \begin{equation}\label{eq:KL-scalar-lower}
        q^{-1}-1+\log q\geq\frac{(q-1)^2}{2\max\{q,1\}^2},
    \end{equation}
    since for $q\geq1$ the left side is $\int_1^q(t-1)t^{-2}\dd t\geq(q-1)^2/(2q^2)$, while for
    $q\leq1$ it is $\int_q^1(1-t)t^{-2}\dd t\geq(1-q)^2/2$. The eigenvalues of $Q$ satisfy
    \begin{equation*}
        q_i-1=-\frac{\Delta t}{1+\Delta t}s_i,
        \qquad
        \max\{q_i,1\}\leq\frac{1+\Delta t\,K}{1+\Delta t},
    \end{equation*}
    with $s_i=1-\mu_i$ the eigenvalues of $S$, so applying \eqref{eq:KL-scalar-lower}
    eigenvalue by eigenvalue gives
    \begin{equation*}
        \KL_C(\rho_{a_{n+1}}\|\rho_{a_n})
        =\frac12\sum_i(q_i^{-1}-1+\log q_i)
        \geq\frac{\Delta t^2}{4(1+\Delta t\,K)^2}\norm S_\Frob^2 .
    \end{equation*}
    Finally $\norm S_\Frob^2=\Tr(C_nRC_nR)=\Tr(C_n[RC_nR])\geq L\Tr(RC_nR)$, because
    $RC_nR\succeq0$ and $C_n\succeq L\Id$. Adding the mean and covariance parts gives
    \begin{equation*}
        \KL(\rho_{a_{n+1}}\|\rho_{a_n})
        \geq
        \frac{L\Delta t^2}{4(1+\Delta t\,K)^2}
        \bigl(\norm g^2+\Tr(RC_nR)\bigr)
        =\frac{L\Delta t^2}{4(1+\Delta t\,K)^2}
        \norm{\grad^{W_2}\calE(a_n)}_{W_2}^2,
    \end{equation*}
    which is the stated bound.
\end{proof}

\subsection{The computable entry gates}

We prove \Cref{prop:entry-residual}. The two gates are the two basic inequalities of
\Cref{sec:global} read in the direction of a stopping test.

\begin{proof}[Proof of \Cref{prop:entry-residual}]
    Write $R(a)=C^{-1}-A(a)$ for the precision residual, so that the Bures--Wasserstein
    residual of the statement is
    \[
        \mathcal R_{\BW}(a)^2
        =\norm{g(a)}^2+\Tr\bigl(R(a)CR(a)\bigr)
        =\norm{\grad^{W_2}\calE(a)}_{W_2}^2 ,
    \]
    which is exactly the left-hand side of the Wasserstein
    Polyak--{\L}ojasiewicz inequality \eqref{eq:W2-PL}. That inequality therefore reads
    $\mathcal R_{\BW}(a)^2\geq2\alpha\,\DeltaE(a)$, and consequently
    $\mathcal R_{\BW}(a_n)^2\leq2\alpha\,\delta_{\loc}$ forces
    $\DeltaE(a_n)\leq\delta_{\loc}$. Every quantity entering
    $\mathcal R_{\BW}(a_n)$ is an expectation under $\rho_{a_n}$ together with $C_n$ and
    $\alpha$, so the test uses no information about $a_\star$.

    For the intrinsic test, suppose $C_n\succeq\ell_n\Id$. The comparison
    \eqref{eq:FR-W2} gives $\Gradnorm(a_n)^2\geq\alpha\ell_n\DeltaE(a_n)$, so
    $\Gradnorm(a_n)^2\leq\alpha\ell_n\delta_{\loc}$ again forces
    $\DeltaE(a_n)\leq\delta_{\loc}$. The first inequality in \eqref{eq:FR-W2} also gives
    $\mathcal R_{\BW}(a_n)^2\leq2\Gradnorm(a_n)^2/\ell_n$, so a triggered Fisher--Rao gate
    triggers the Bures--Wasserstein gate as well, which is the asserted comparison between
    the two tests.

    Finally, both deterministic schemes decrease the energy at the stepsizes of
    \Cref{thm:glob-Riem,thm:glob-KL}, and the flow decreases it by
    \eqref{eq:dissipation}. Hence $\DeltaE(a_k)\leq\DeltaE(a_n)\leq\delta_{\loc}$ for every
    $k\geq n$, so once either test triggers the iterate remains in the certified sublevel
    set permanently.
\end{proof}

% SOURCES:
%   Proof of prop:entry-residual assembled from the two basic inequalities of Section 2
%   (eq:W2-PL, eq:FR-W2, both proved there) plus the energy-descent property of the two
%   schemes; cf. improved-global-local.tex l.4088-4176 (prop:entry-residual).
%   improved-global-local.tex l.1850-1877 (lem:scalar-retraction) -> lem:scalar-retraction
%   improved-global-local.tex l.2115-2162 (lem:scalar-KL), identical to
%     improved-global-local-stoch.tex l.4030-4058 -> lem:scalar-KL
%   improved-global-local.tex l.1879-1991 (lem:Riem-descent, retraction smoothness for
%     arbitrary tangent directions and the deterministic descent step) -> proof of lem:descent-Riem
%   improved-global-local-stoch.tex l.3994-4025 (lem:KL-argmin) -> lem:KL-argmin
%   natural-gradient.tex l.564-644 (prop:conv-smooth-D, the author's own square-root-factor
%     argument with the G_0 >= 0 step, base point renamed a_0 -> a_n) -> prop:smooth-Bregman
%   improved-global-local.tex l.2164-2285 (prop:KL-onestep) and
%     improved-global-local-stoch.tex l.4063-4118 (lem:KL-onestep-app, the self-contained
%     version) -> proof of prop:onestep-KL
